\section{Related work}\label{sec:related-work}

\Flux's claim in this paper is architectural, not comparative marketing: a network in which an
autonomous agent can rent capacity, pay for it, and leave, without an account, a payment card, or a
human in the loop, using an operator set that is capital-bonded and enumerable rather than merely
reputational. Positioning that claim requires comparing mechanisms, not brands. This section takes
each mechanism the paper specifies elsewhere — a decentralized resource market, a consensus protocol,
a bridge, an attestation root, a shielded payment pool, a delegation primitive, and a machine-facing
interface — and states the prior art for that mechanism specifically, what \Flux inherits from it, and
where \Flux's design departs. Claims about other systems describe only what those systems have
published, with no maturity judgement about their roadmaps; claims about \Flux itself carry the same
maturity and provenance discipline as the rest of this paper.

\subsection{Decentralized compute markets}

Akash \citep{akash} and Golem \citep{golem} are open compute marketplaces: a tenant posts a resource
request, providers bid or advertise matching capacity, and a match is made without the marketplace
itself verifying that a provider's declared hardware exists or performs as advertised; both settle
payment in the network's own token, escrowed from the tenant and released to the winning provider over
the deployment's life. Render \citep{render} and io.net \citep{ionet} apply the same open-bidding
structure to a single resource class — GPU rendering and GPU compute for machine-learning workloads,
respectively — aggregating capacity from an unbonded, self-reported pool of node operators. Aethir
\citep{aethir} targets the same GPU-cloud niche with an enterprise framing. In all five, the trust a
tenant extends to a provider is reputational: nothing in the protocol requires a provider to lock
capital against non-delivery, and nothing verifies a provider's declared specification beyond the
provider's own attestation of it.

\Flux departs from this class in one specific way: every \FluxNode operator posts a tier-scaled
collateral that consensus checks against the registration transaction before a node is eligible to run
anything \shipped\ \srccode{fluxd/\allowbreak{}src/\allowbreak{}coins.cpp}{630--686}, and the resulting registry is directly
enumerable from chain state — 6{,}489 records at the most recent measurement \shipped\
\srcmeasured{api.runonflux.io/\allowbreak{}daemon/\allowbreak{}listzelnodes}{2026-09-03}. Two properties follow that an open,
self-reported marketplace cannot offer by construction. First, the same registry that gates eligibility
is what \PoN's block-production sortition ranks and rotates \shipped\
\srccode{fluxd/\allowbreak{}src/\allowbreak{}pon/\allowbreak{}pon-minter.cpp}{78--104}, and it is the substrate §7's fairness argument for
revenue-funded compensation is built on \planned\ \srcintent{§7}. Second, an enumerable operator set is
a tractable target for a future key migration in a way an open, unbonded pool cannot be: collateral is
already a registered set on \Flux, so a migration that starts there is a narrower problem than a
general UTXO migration \vision\ \srcintent{feasibility assessment}. This is also what makes the resulting
eviction quorum Sybil-resistant by construction \citep{douceur2002sybil} — a property §23 relies on and
this section returns to below. Neither property is free: collateral is capital locked against no
guaranteed return, and an enumerable registry is also a measurably concentrated one — the largest single
collateral-weighted identity in the registry holds 10.71\% of voting power across 237 Stratus nodes
under one key \shipped\ \srcmeasured{api.runonflux.io/\allowbreak{}daemon/\allowbreak{}listzelnodes}{2026-09-03}, a fact §23 uses
directly.

\subsection{Hyperscalers}

Hyperscale clouds occupy an adjacent but distinct point in the design space, and there is no hostility
in stating the difference plainly. They price by metered consumption rather than by reservation, offer
deep managed-service catalogs, and back regional deployments with contractual latency and availability
guarantees; none of that is disputed here \citep{wust2018blockchain}. What \Flux offers that a
hyperscaler does not is orthogonal rather than superior: no account, no payment card, and no human
approval step stand between a payment and a running deployment — already true, concretely, of
\CumulusVPN's entitlement path, where a WireGuard public key is itself the credential and no
server-issued account exists \shipped\ \srccode{cumulusvpn/\allowbreak{}gateway/\allowbreak{}internal/\allowbreak{}entitle/\allowbreak{}entitle.go}{218--287}
— and the operator fulfilling the deployment is drawn from the capital-bonded, enumerable set described
above rather than an opaque corporate fleet. Concretely, \Flux prices a reservation: a deployment's
price is a function $\price(\resvec,\ninst,\dur)$ of a resource vector, an instance count and a duration fixed at request time,
not a metered function of realized usage \shipped\
\srccode{flux/\allowbreak{}ZelBack/\allowbreak{}src/\allowbreak{}services/\allowbreak{}utils/\allowbreak{}appUtilities.js}{89--134}
\srcmeasured{api.runonflux.io/\allowbreak{}apps/\allowbreak{}deploymentinformation}{2026-09-03} (§16). That is the architectural
difference that matters more than any feature count: a reservation price is settled once, at deployment
time, and the only subsequent dispute surface is whether the reserved resources were delivered, not how
much of them was consumed.

\subsection{Storage and content networks}

Filecoin \citep{filecoin2017} and IPFS \citep{benet2014ipfs} are the closest prior art for hosting
rather than compute. Filecoin's storage deals are prepaid by the client and backed by provider-posted
collateral that the protocol can slash if the provider fails to prove continued possession of the
stored data; the remedy for non-performance is financial, enforced against the provider's own stake.
\Flux's stated remedy for a host that fails to deliver a paid-for deployment is different in kind:
replacement rather than refund — the network moves the workload to another eligible host rather than
returning payment to the tenant \planned\ \srcintent{§16}. IPFS carries no equivalent mechanism at
either layer: content addressing guarantees that a hash resolves to the bytes it names, but the network
has no built-in remedy for a bad actor and no protocol-level way to remove content once published —
persistence is a function of who chooses to pin it, not of any consensus process. \Flux's proposed
content-reporting quorum (§23) is the opposite kind of mechanism: a collateral-weighted vote to evict a
specific application from the fleet \planned\ \srcintent{§23}. Placed side by side, the three positions
are distinct: Filecoin slashes a provider's stake for failing to keep data available; IPFS has no
mechanism to make content unavailable at all; \Flux's quorum is built expressly to make a specific
application unavailable network-wide, which is not the same problem as either.

\subsection{Consensus}

\PoN replaces Bitcoin-style proof of work \citep{nakamoto2008bitcoin} with collateral-gated sortition
over the \FluxNode registry described above: eligibility to produce the next block is a function of
tier and collateral rather than of hashing power. The precise point a reader should not have to infer
is in the fork-choice rule. Every \PoN block contributes a fixed chainwork value of $2^{40}$ regardless
of its actual difficulty target \shipped\ \srccode{fluxd/\allowbreak{}src/\allowbreak{}pow.cpp}{284--290}, so fork choice is not a
race for the chain with the most accumulated work in the sense Nakamoto consensus assumes
\citep{nakamoto2008bitcoin} and the security analyses built on that assumption use
\citep{gervais2016security,eyal2014majority} — it is a block-count race, resolved by which chain
extended first and furthest. That places \PoN's security argument closer to proof-of-stake sortition
schemes such as Algorand \citep{gilad2017algorand}, which likewise select a block producer by a
verifiable, stake-weighted draw rather than by cumulative work, than to Nakamoto consensus or to
chain-based proof-of-stake protocols such as Ouroboros \citep{kiayias2017ouroboros} that still
accumulate a work-like quantity along the chain \srcinferred. Two differences from Algorand are worth
stating rather than eliding: \PoN's eligibility draw is public and registry-based rather than a private
per-round cryptographic sortition, and, so far as the material available to this paper shows, \PoN has
no finality gadget layered over its fork-choice rule, unlike systems that pair a longest-chain-style
base protocol with an explicit finality overlay \citep{buterin2017casper} \srcinferred.
\citep{bonneau2015sok} frames exactly this kind of comparison — what a consensus protocol's security
argument actually rests on, rather than which family it is conventionally filed under — and is the
right lens for reading \PoN this way.

\subsection{Bridging}

Cross-chain interoperability designs separate broadly into light-client verification, trusted or
federated attestation, and atomic exchange of assets that already exist on both chains
\citep{zamyatin2021sok}. IBC \citep{cosmosibc} is the clearest instance of the first category: a
receiving chain runs a light client of the sending chain, verifying block headers and validator
signatures directly, so no third party is trusted beyond the sending chain's own consensus. XCLAIM
\citep{zamyatin2019xclaim} generalizes the idea to chains, like Bitcoin, that cannot be light-client
verified cheaply on the receiving side, using over-collateralized vaults and SPV-style proofs instead.
CCTP \citep{cctp} sits in the second category: an attestation service observes a burn on the source
chain and signs a message authorizing a mint on the destination chain, substituting a designated
attester for a light client. Atomic swaps \citep{herlihy2018atomic} sit outside both categories: they
exchange assets that already exist on both chains using hashed time-locks, but they cannot
\emph{issue} an asset on a chain where it does not yet exist.

Neither of the first two categories is presently constructible for \Flux, and the reason is the header,
not the will. A light client needs three things \PoN's header does not give it. First, per the bridge
design's own analysis, the header carries no commitment to the \FluxNode registry, so a verifier cannot
check that a block's signer was an eligible operator without holding the full registry itself \planned\
\srcintent{mechanism specification}. Second, chainwork is fixed at $2^{40}$ per block regardless of
difficulty \shipped\ \srccode{fluxd/\allowbreak{}src/\allowbreak{}pow.cpp}{284--290}, so there is no cumulative-work quantity a
light client could use to prefer one candidate chain over another the way an SPV client does. Third, a
2-of-4 emergency key set can produce a valid block outside \PoN's normal eligibility rule, with no time
or stall-detection gate on when it may act \shipped\ \srccode{fluxd/\allowbreak{}src/\allowbreak{}emergencyblock.cpp}{183--191}, so no
light-client rule derived from \PoN eligibility alone would be sound even if the first two problems were
solved. The consequence is concrete: the most trust-minimized bridge architecture available to
IBC-style systems is closed to \Flux until the header changes to commit to the registry, and that is a
consensus change, not a bridge-layer one, and a candidate for the same network upgrade that would carry
collateral maturity \planned\ \srcintent{§8}.

Atomic swaps do not fill this gap either, for a reason independent of the header: a swap exchanges FLUX
that already exists on \Flux for an asset that already exists on the destination chain. It cannot mint
FLUX on a chain where none yet exists, which is exactly what parallel-asset issuance requires \planned\
\srcintent{mechanism specification}. Swaps are the right primitive for \FusionX's exchange product, where
both assets already exist on both sides; they are not a candidate for the bridge itself. \Flux's
remaining option, and the one §22 specifies, is canonical mint-and-burn under a bonded, capped
custodian — closer to CCTP's trusted-attester model than to IBC's light client, with per-chain caps
$\bcap{c}$ and reserve slack $\bslack = \breserve - \sum_c \bminted{c}$ (§22) doing the work a light
client would otherwise do \planned\ \srcintent{§22}.

\subsection{Attestation and bonded assertions}

ArcaneOS's measured boot and Secure-Boot platform-key pinning \shipped\ \srccode{fluxsas/\allowbreak{}src/\allowbreak{}fes/\allowbreak{}fes.cpp}{70--81},
its dm-verity root-hash verification \shipped\ \srcprivate{\SAS{} watchdog}, and its
LUKS-backed disk encryption \shipped\ \srcprivate{\SAS{} disk encryption} do real work
against a remote attacker without any special hardware trust anchor. The question this section
positions is not whether that machinery works but what the resulting attestation is rooted in. One line
of prior art roots attestation in silicon: Intel SGX \citep{costan2016sgx} ties an attestation to a
hardware-manufacturer-signed key that a remote verifier can check without trusting the party running
the enclave. A second line roots it in capital: EigenLayer \citep{eigenlayer} lets operators post
slashable stake against off-chain assertions, so an assertion is trusted because a bond would be lost
if it is false, not because of where it was computed. \Flux's planned attestation network is a version
of the second model — a bond rather than a chip is what would secure the claim that a machine booted a
particular image \planned\ \srcroadmap{new benchmark tool, signed node attestation,
timeline.html:295--296}.

The honest state of the current system is that neither model is yet in place. The \fluxsas signing key
for the default, \texttt{cdn}, \texttt{mesh} and \texttt{app} attestation purposes is derived from
static salt and domain values compiled identically into every binary; none of its inputs depends on TPM
state, Secure Boot state, or any other per-machine value \shipped\
\srcprivate{\SAS{} API server}. What that attestation demonstrates is that the signer
holds a value shipped in every copy of the software, not that a particular machine booted a particular
image. The trust root behind "this node booted something \Flux considers legitimate" is today entirely
\Flux-operated, composed of the platform key above and a pair of \Flux-operated web services publishing
allowed dm-verity root hashes that the watchdog polls roughly daily and treats as secure if unreachable
\shipped\ \srcprivate{\SAS{} watchdog}. Neither SGX's hardware root nor EigenLayer's
economic root is present yet; the \texttt{purpose=\allowbreak{}mesh} attestation type already compiled into
\fluxsas with no consumer \shipped\ \srcprivate{\SAS{} API server} is the stated seed of
the bonded design that would supply one \planned\ \srcintent{§14}.

\subsection{Privacy}

\Flux's shielded pool is Sapling — encrypted notes, a note commitment tree $\tree$, nullifiers $\nf$,
and Groth16 proofs \citep{groth2016size} verifying spend and output statements over $\zaddr$-typed
value — specified by the Zcash protocol \citep{zcashprotocol} and tracing to the original Zerocash
construction \citep{bensasson2014zerocash}, and it is live on-chain today \shipped\
\srccode{fluxd/\allowbreak{}src/\allowbreak{}main.cpp}{1398--1408}. The trusted setup behind those proofs was produced by a
multi-party ceremony in the random-beacon model \citep{bowe2017mpc}; a newer construction removing the
need for any such ceremony exists in the literature \citep{bowe2019halo} but has not been ported into
\Flux's shielded pool, which remains on Sapling and Groth16 \shipped\ \srccode{fluxd/\allowbreak{}src/\allowbreak{}main.cpp}{1450--1468}.

\textbf{The comparison changes direction going forward.} Both pools are to be retired
(\S\ref{sec:shielded-retirement}), so \Flux{} is moving \emph{away} from the Zcash lineage on this
axis while the systems it is compared against move toward stronger privacy. The resulting position is
unusual and should be stated as such rather than blurred: on settlement privacy \Flux{} will be
comparable to Bitcoin rather than to Zcash, and weaker than every privacy-focused chain in this
survey. What it gains is a consensus layer with no proving system, no trusted-setup inheritance and
no unenforced turnstile --- a smaller and more auditable object, which is the trade the team has
chosen \srcintent{design intake}.
None of this machinery is \Flux's invention, and this paper does not claim otherwise: the cryptography,
the setup, and the protocol are inherited in full.

Where \Flux does make an original decision is consensus policy, not cryptography: since height
835{,}554, the same consensus path cited above rejects any transaction that mixes a transparent input
with a shielded output or joinsplit, closing the pool to new value from the transparent side \shipped\
\srccode{fluxd/\allowbreak{}src/\allowbreak{}main.cpp}{1398--1408}. Value already inside the pool can still move
$\zaddr\to\zaddr$ or leave $\zaddr\to\taddr$, but nothing enters; the pool's value has been
monotonically non-increasing for approximately 5.4 years, holding 77{,}188.76 FLUX in Sapling and 19{,}291.18 FLUX in
Sprout at last measurement \shipped\ \srcmeasured{fluxd chain tip 2{,}917{,}115}{2026-09-03}. That
single fact reorders how privacy commitments elsewhere should be read: reopening the pool would have
been a consensus change and the prerequisite for every downstream privacy item, not a
wallet-integration task --- and it is now cancelled in favour of retirement (\S\ref{sec:shielded-retirement}).

Payment channels are the other well-known route to transactional privacy — Lightning
\citep{poon2016lightning} keeps individual payments off-chain entirely, trading on-chain settlement for
a channel-liquidity trust assumption. \Flux does not take this route for deployment payment; §16 works
over on-chain transparent settlement and §21 was specified over on-chain shielded settlement \srcintent{§16}, so a channel network's
liquidity and routing assumptions are not this paper's problem, but neither is its privacy benefit
available for free. The specific problem §21 identifies is narrower and, so far as this corpus
establishes, open: proving that a shielded payment funds a \emph{particular} deployment without
revealing which one is hard precisely because an application's address is itself a workload identifier,
so a spend statement naming that address defeats the purpose of hiding it \vision\
\srcintent{feasibility assessment}. This has the shape of a blind-credential problem — proving membership
in, or payment toward, a class without revealing which member —
and the relevant literature is well developed: blind signatures \citep{chaum1982blind},
anonymous credentials \citep{camenisch2001credentials}, compact e-cash
\citep{camenisch2005ecash}, anonymous authorization tokens \citep{privacypass}, and
zero-knowledge group membership with per-epoch nullifiers \citep{semaphore}. What none of it
supplies directly is the combination \S\ref{sec:private-payment} needs: a public obligation
checkable by thousands of independent verifiers holding no shared secret, settled from a pool that
already exists but cannot be entered, at a demand level where the anonymity set rather than the
cryptography is the binding constraint.

\subsection{Delegated authority}\label{sec:rw-delegation}

ERC-4337 \citep{erc4337} answers the delegation problem by putting a smart-contract account between a
session key and the assets it can move: a paymaster or account contract can enforce a spending policy
on-chain, because the account itself is a program the chain executes. That answer is not available on
\Flux's L1, which is a UTXO chain: script validation is a pure function of a spending transaction and
its own inputs, with no access to spending history, so no script-level rule can enforce a rate limit or
a cumulative cap on a delegated key (§17; \srcintent{feasibility assessment}). Macaroons
\citep{birgisson2014macaroons} fit the UTXO setting better than account abstraction does, because they
push attenuation into the credential itself rather than into on-chain account state: a macaroon's
caveats are checked by whoever verifies the credential, not by a ledger, which matches how deployment
specifications are already validated off-chain against owner signatures on every node (§19;
\srcintent{feasibility assessment}). That correspondence is the basis for §17's specification of bounded
deployment authority \planned\ \srcintent{§17}. It is worth stating precisely what already exists and
what does not: a scoped, revocable, rate-limited bearer token already governs \emph{actions} — start,
stop, deploy, renew — through the FluxEdge API (§17; \srcintent{feasibility assessment}), and an
MCP-exposed tool set can already sign and submit a deployment end to end on a custodial path with no
bound on \emph{value} at all \shipped\
\srccode{deploy-with-git-frontend/\allowbreak{}server/\allowbreak{}mcp/\allowbreak{}createServer.js}{174--266}. §17 is therefore not inventing
delegation from nothing; it specifies the value bound that today's action-scoped delegation does not
have.

\subsection{Agent interfaces}

The Model Context Protocol \citep{mcp} is an emerging, vendor-published convention for exposing tools
to a language-model agent over a uniform interface, rather than a bespoke API per application. \Flux
already exposes deployment operations this way: \texttt{deploy\_app}, \texttt{update\_app} and
\texttt{renew\_app} are registered as MCP tools today \shipped\
\srccode{deploy-with-git-frontend/\allowbreak{}server/\allowbreak{}mcp/\allowbreak{}createServer.js}{174--266}. §18 gives the fuller listing
and states what the interface does not yet bound, which is the same value-bound gap \S\ref{sec:rw-delegation} describes.

\subsection{What this paper adds}

Set against this prior art, three results in this paper are, to the corpus available here, original
rather than inherited, and each is narrow. First, the fee-pool drip construction in §7 — paying the
entire \PoN rotation from a shared pool $\pool$ at a fixed per-block rate $\pool/\drip$ under \PNR — is
offered as a specific defence against a timing attack that a naive next-block payout admits; the
mechanism composes an existing rotation with a new payment rule, and the fairness argument is this
paper's, not inherited from any cited system \planned\ \srcintent{§7}. Second, §23's finding is a
measurement, not a mechanism: collateral-weighted moderation is Sybil-resistant by construction
\citep{douceur2002sybil}, and, at the measured distribution — a Gini coefficient of 0.8375, with four
owner identities (three under linkage) alone clearing the 25\% eviction threshold — it is not censorship-resistant,
and no reweighting inside the design space closes that gap \shipped\
\srcmeasured{api.runonflux.io/\allowbreak{}daemon/\allowbreak{}listzelnodes}{2026-09-03}. Third, §22 and §4 together identify a
foreclosure rather than a new construction: a header format chosen for \PoN's own consensus reasons —
fixed chainwork, no registry commitment, an ungated emergency path — rules out the light-client bridge
architecture standard for account-based chains, a cost of that consensus design not stated elsewhere.
None of the shielded pool, the Sapling proof system, the collateral-sortition idea itself, MCP, or
macaroons is claimed as this paper's contribution; each is prior art this paper positions \Flux
against, mechanism by mechanism.

%% Drain this section's float queue before the next begins.
\FloatBarrier
